\documentclass[dvipdfmx]{jsarticle}
%dvipdfmxを必ず指定する。
\usepackage{mypreamble}
%%%もろもろの定義が長かったため，mypreamble.styにまとめています。
\usepackage{ascolorbox}
\title{tcolorboxのお誘い}
\author{ぱるち}
\date{2020年6月15日}

\begin{document}

\maketitle

なお，今回のpdfを作成するにあたり，和文フォントを2種類より多く使用するパッケージを導入しています。
\begin{itemize}
    \item \verb|\usepackage[deluxe]{otf}|
    \item \verb|\usepackage[noto-otc,unicode]{pxchfon}|
\end{itemize}

\begin{marker}
\verb|\usepackage[noto-otc,unicode]{pxchfon}|については，texlive2019やtexlive2020で正常に動作します。もしもpdfのレンダリングがうまくいかない場合は，このパッケージの読み込みを取り消してください。
\end{marker}

\section{基本的な例}

まずは最も基本的な例。

\begin{tcolorbox}
Hello, \TeX !!
\end{tcolorbox}

次は色を付けたりタイトルをつけたり。

\begin{tcolorbox}[title=これがタイトルで2番目の例,
fonttitle=\gtfamily\sffamily\bfseries,
colframe=blue,colback=blue!3!]
にほんごでも全然かまわないのです。
\end{tcolorbox}

本文中でも新しい環境を定義することは可能。煩雑になるので推奨しませんが。。

%%% boxmine環境の定義
\newtcolorbox{boxmine}[2][]{colbacktitle=red!10!white,
colback=blue!10!white,coltitle=red!70!black,
title={#2},fonttitle=\bfseries\gtfamily,#1}
%%%

\begin{boxmine}{ここにTitleを書くのだ}
吾輩は猫である。名前はまだない。\\
どこで生れたか頓（とん）と見当がつかぬ。何でも薄暗いじめじめした所でニャーニャー泣いていた事だけは記憶している。吾輩はここで始めて人間というものを見た。しかもあとで聞くとそれは書生という人間中で一番獰悪（どうあく）な種族であったそうだ。
\end{boxmine}

\begin{boxmine}[fontupper=\gtfamily\Large,
before upper=【大きくしてみた。】]{さらにオプションを加えるのだ}
吾輩は猫である。名前はまだない。\\
どこで生れたか頓（とん）と見当がつかぬ。何でも薄暗いじめじめした所でニャーニャー泣いていた事だけは記憶している。吾輩はここで始めて人間というものを見た。しかもあとで聞くとそれは書生という人間中で一番獰悪（どうあく）な種族であったそうだ。
\end{boxmine}


\section{いろいろな枠を作って見よう}

この節以降の枠については，定義を\verb|mypreamble.sty|に書いています。

第1節と第2節で紹介のレベルにギャップがあるかもしれませんね。そのうちギャップを埋めたいなとも考えています。

\subsection{問題集の例題風}

自作の枠です。reidai環境の定義は\verb|mypreamble.sty|の12行目から。

\begin{reidai}
$a$，$b$は実数で，$f(x)=x^2+ax+b$とする。$\alpha$，$\beta$を2次方程式$f(x)=0$の異なる2解とする。$\alpha^2$，$\beta^2$がまた$f(x)=0$の異なる2つの実数解であるとき，$a$，$b$の値を求めよ。
\end{reidai}

下部に解答を書くことも可能です。

\begin{reidai}
次の問題に答えなさい。
\begin{enumerate}
    \item 8人を2つの組に分ける方法は何通りあるか。
    \item 6人を3つの部屋A，B，Cに入れる方法は何通りあるか。ただし各部屋に少なくとも1人は入るものとする。
\end{enumerate}

\tcblower

区別があるかどうかを正しく考えます。
\begin{enumerate}
    \item なんだかんだで127通り
    \item なんだかんだで540通り
\end{enumerate}

\end{reidai}

999番までは問題ない出力が得られると思います。1000番以降になるとはみ出してしまいます。この辺りは余白を設定すれば十分ですが。

\setcounter{reidaibangou}{998}
\begin{reidai}
$a$，$b$は実数で，$f(x)=x^2+ax+b$とする。$\alpha$，$\beta$を2次方程式$f(x)=0$の異なる2解とする。$\alpha^2$，$\beta^2$がまた$f(x)=0$の異なる2つの実数解であるとき，$a$，$b$の値を求めよ。
\end{reidai}

\begin{reidai}
$a$，$b$は実数で，$f(x)=x^2+ax+b$とする。$\alpha$，$\beta$を2次方程式$f(x)=0$の異なる2解とする。$\alpha^2$，$\beta^2$がまた$f(x)=0$の異なる2つの実数解であるとき，$a$，$b$の値を求めよ。
\end{reidai}
\verb|\ref{}|を使ってカウンターを参照する機能は実装していませんが，\verb|\tcbuselibrary{theorems}|の機能を使えば能うハズ。

\subsection{教科書風に公式を書く}

個人的には結構きれいな出来になったんじゃないかなあと思っています。

\begin{kousiki}{中間値の定理}
区間$[\alpha,\beta]$で連続な関数$f(x)$について，
$f(\alpha)$と$f(\beta)$の間にある任意の実数$c$に対して，
ある実数$k\in (\alpha,\beta)$を\[
f(k)=c
\]
を満たすようにとることが出来る。
\end{kousiki}

\subsection{定理環境}
最近僕がtexで文書を書く際は，\verb|asmthm.sty|で定理環境を定義するよりも，\verb|tcolorbox.sty|を用いた定義が多いです。
\begin{mytheo}{中間値の定理}{chukan}
区間$[\alpha,\beta]$で連続な関数$f(x)$について，
$f(\alpha)$と$f(\beta)$の間にある任意の実数$c$に対して，
ある実数$k\in (\alpha,\beta)$を，$f(k)=c$を
満たすようにとることが出来る。
\end{mytheo}
例えば\ref{tha:chukan}とすれば，上の中間値の定理の番号が出力されます。もちろんハイパーリンクを付けることも，\verb|hyperref.sty|を読み込めば良いです。

記事では紹介しませんでしたが，定理と命題などを連番にすることも可能です。


\begin{myprop}{方程式の実数解の存在}{}
区間$[\alpha,\beta]$で連続な関数$f(x)$について，
$f(\alpha)f(\beta)<0$ならば，方程式$f(x)=0$は$\alpha<x<\beta$の範囲に少なくとも1つの実数解をもつ。
\end{myprop}

\section{tcolorboxからの引用枠}
単純にすごいです。

\begin{marker}
私は場合もっともこの反駁学というのの時へ閉じたませ。\\
どうしても今に反抗人はまあそうした研究でなくだけのしからみるたには矛盾さならだて、当然にもしないたたです。腹の中に行きた事も初めて十月にともかくうますない。
\end{marker}

\section{ascolorbox.styを使ってみる}

ここでは鉄緑会の方々によって作られた\verb|ascolorbox.sty|を用いた例を紹介します。githubからインストール可能です。URLは\url{https://github.com/Yasunari/ascolorbox}

\begin{practicebox}{中間値の定理}
地球の赤道上にある点Pを置き，その対蹠地をQとする。ここでQは赤道上にあると仮定する。このとき点Pでの気温と点Qの気温が同じになるように，点Pを上手にとることができる。その理由を説明しなさい。
\end{practicebox}


\begin{ascolorbox8}{コラム}
私は場合もっともこの反駁学というのの時へ閉じたませ。\\
どうしても今に反抗人はまあそうした研究でなくだけのしからみるたには矛盾さならだて、当然にもしないたたです。腹の中に行きた事も初めて十月にともかくうますない。
\end{ascolorbox8}

\end{document}
